\documentclass[12pt]{article}
\usepackage[utf8]{inputenc}
\usepackage[T1]{fontenc}
\usepackage{hyperref}
\title{Yeni Language Resources for K-12 Curriculum Development}
\author{Compiled by Claude Sonnet 5.5 (claude-sonnet-5-5) for Andr\'as Kornai}
\date{}
\begin{document}
\maketitle

\section{Introduction}
Yeni (ISO 639-3 \texttt{yei}, Glottocode yeni1253) is an extinct language of Cameroon. It was once spoken around Djeni Mountain in the Nyalang area. According to Wikipedia, only a song remembered by some Sandani speakers survives. Bruce Connell reported it in 1995. From the song's vocabulary he judged it closely related to the Mambiloid languages Cambap, Njerep and Kasabe. Its classification is uncertain.

Caution: the ``Yei'' entry at Global Recordings Network has ISO code \texttt{jei}. It is a different language and must not be confused with Yeni.

\section{Primary Educational Resources}
None exist. There are no speakers and no teaching materials.

\section{Linguistic and Reference Materials}
\begin{itemize}
\item Wikipedia article on Yeni: \url{https://en.wikipedia.org/wiki/Yeni_language}
\item Glottolog languoid page for Yeni (yeni1253): \url{https://glottolog.org/resource/languoid/id/yeni1253.bigmap.html}
\end{itemize}

\section{Religious and Community Resources}
None were found.

\section{Digital and Audio Resources}
None were found. The one surviving song was not located as a recording.

\section{Curriculum Development Context}
The evidence is too thin for a language curriculum. A heritage unit could cover the language's history and its place in Cameroon's linguistic diversity. It could also use the related Mambiloid languages as context, without presenting reconstructed Yeni as attested.
\end{document}